\section{Prescribing the slow mode: a two-dimensional feasibility test}\label{sec:realizability}
The thermal classification concerns an eigenvector of the full generator.
There is also a finite inverse test for whether a proposed direction can occur
while every original triplet test fails. It reduces the rate-design problem,
including spectral ordering, to the intersection of a disk and eight half-planes.

Fix the bath probabilities $q$ and a nonconstant vector $w$ with $q^\top w=0$.
Let $D=\operatorname{diag}(q)$, and choose any full-rank $4\times2$ matrix $U$
whose columns span
\[
\{f:q^\top f=0,\ (Dw)^\top f=0\}.
\]
For a symmetric $2\times2$ matrix $B$, set
\[
M_0(B)=-DU B U^\top.
\]
This auxiliary matrix need not have nonnegative transition rates; it has
zero column sums, detailed balance, and null left vectors $\mathbf1,w$.

\begin{theorem}[Exact inverse feasibility]\label{thm:inversefeasibility}
There exists a strictly positive reversible generator for which $w$ is the
simple slow left mode and every original triplet predicate fails if and only
if some positive definite $B$ satisfies the eight linear inequalities
$F_{ijk}(M_0(B))\leq0$ and $G_{ijk}(M_0(B))\leq0$.
Equivalently, normalize $\operatorname{tr}B=1$ and write
\[
B=\frac12\begin{pmatrix}1+x&y\\y&1-x\end{pmatrix}.
\]
The exact feasibility condition is that the resulting eight closed
half-planes intersect $x^2+y^2<1$.
\end{theorem}
\begin{proof}
If $B$ is positive definite, $D^{-1/2}M_0D^{1/2}$ is negative semidefinite,
with exactly the two-dimensional kernel spanned by $\sqrt q$ and
$D^{1/2}w$. Add a reset generator:
\[
M=M_0+\rho(q\mathbf1^\top-I).
\]
For sufficiently large $\rho>0$, all off-diagonal entries are strictly
positive. The nonstationary eigenvalues are $-\rho,-\rho-\nu_1,
-\rho-\nu_2$ with $\nu_1,\nu_2>0$, so $w$ is the simple slow mode.
The reset contributes zero to every $F$ and $G$, hence the inequalities
persist. Since the physical rates are now positive, the branch-equivalence
lemma identifies these with failure of the original predicates.

Conversely, start with any such physical generator, with simple slow
eigenvalue $\lambda_2<0$. Put
$M_0=M+\lambda_2(q\mathbf1^\top-I)$. Its symmetrization is zero on the
stationary and slow directions, and strictly negative on their two-dimensional
orthogonal complement. Therefore $M_0=-DU B U^\top$ for a positive definite
$B$. Reset invariance gives the eight inequalities. Positive scaling of $B$
preserves them, so trace normalization loses no solution. A symmetric trace-one
$2\times2$ matrix is positive definite exactly when $x^2+y^2<1$.
\end{proof}

For rational $q,w$, choose $U$ rational by Gaussian elimination. All eight
half-planes then have rational coefficients. Exact feasibility is particularly
simple: find the point of their intersection closest to the origin, and compare
its squared norm with one. The closest point, when the polyhedron is nonempty,
is either the origin, the orthogonal projection of the origin onto one boundary
line, or the intersection of two nonparallel boundary lines. Enumerate these
finitely many candidates and discard those violating any half-plane. Each
remaining squared norm is rational. A minimum below one certifies realizability;
a minimum equal to or above one, or no feasible candidate, excludes a simple
slow realization. Equality must not be rounded into an admissible positive
definite matrix.

This gives a complete checkable inverse characterization at fixed rational
bath weights and rational mode direction. It uses established reset and
semidefinite geometry, but the half-planes enforce the exact original triplet
premise. It is not a claim that eigenvectors depend convexly on rates. It also
explains why prescribing a thermal shape can make adversarial searches much
more efficient than sampling rate matrices: spectral ordering is enforced by
positive definiteness of only a $2\times2$ matrix.

\paragraph{The necessary slow-mode inequalities are not sufficient.}
The structural restriction proved in Section~\ref{sec:fulltripletbound}
does not replace this inverse test. For an exact obstruction, use
unnormalized bath weights $\widehat q=(5,4,2,1)$ and
\[
w=(162,-18,71,-880).
\]
Its bath mean is zero, its three Helmert coefficients are $(80,2,80)>0$,
and $w_1>w_2$, $w_1+w_2-2w_3=2>0$, $w_4<0$. Nevertheless no positive
full-triplet generator can have this simple slow direction. To see this,
use $\widehat D=\operatorname{diag}(\widehat q)$ and
\[
U=\begin{pmatrix}
-89/450&431/450\\ -91/360&-521/360\\1&0\\0&1
\end{pmatrix}.
\]
For the signed core $M_0=-\widehat DUBU^\top$, exact expansion gives
\[
\frac{482571629}{11676414}F_{123}(M_0)
+\frac{496382581}{11676414}F_{124}(M_0)
+\frac{16620682809}{77842760}G_{234}(M_0)
=\operatorname{tr}B.
\]
All coefficients are positive. The triplet constraints would imply
$\operatorname{tr}B\leq0$, contradicting $B\succ0$. Using unnormalized
weights only rescales the fast matrix by a positive constant and does not
change feasibility. Reset addition leaves these expressions unchanged.
This is a Farkas-type separating certificate, checked by
\texttt{extended\_mode\_obstruction\_certificate.py}; it explains why the
new structural inequalities remain necessary rather than a complete
realizability characterization.
